% Part: first-order-logic
% Chapter: tableaux
% Section: identity

\documentclass[../../../include/open-logic-section]{subfiles}

\begin{document}

\olfileid[pt]{fol}{tab}{ide}

\olsection{\usetoken{P}{tableau} com \usetoken{S}{identity}}

!!^{tableau}s com !!{identity} exigem regras de inferência adicionais.
As regras para $\eq$ são as seguintes ($t$, $t_1$ e $t_2$ são termos fechados):

\begin{defish}
\AxiomC{}
\RightLabel{$\eq$}
\UnaryInfC{\sFmla{\True}{\eq[t][t]}}
\DisplayProof
\hfill
\AxiomC{\sFmla{\True}{\eq[t_1][t_2]}}
\noLine
\UnaryInfC{\sFmla{\True}{!A(t_1)}}
\RightLabel{$\TRule{\True}{\eq}$}
\UnaryInfC{\sFmla{\True}{!A(t_2)}}
\DisplayProof
\hfill
\AxiomC{\sFmla{\True}{\eq[t_1][t_2]}}
\noLine
\UnaryInfC{\sFmla{\False}{!A(t_1)}}
\RightLabel{$\TRule{\False}{\eq}$}
\UnaryInfC{\sFmla{\False}{!A(t_2)}}
\DisplayProof
\end{defish}
Observe que, diferentemente de todas as outras regras, $\TRule{\True}{\eq}$ e
$\TRule{\False}{\eq}$ exigem que \emph{duas} !!{formula}s sinalizadas
já apareçam no ramo, a saber, tanto $\sFmla{\True}{\eq[t_1][t_2]}$
quanto $\sFmla{S}{!A(t_1)}$.

\begin{ex}
Se $s$ e $t$ são termos fechados, então $\eq[s][t], !A(s)
\Proves !A(t)$:
\begin{oltableau}
  [\sFmla{\False}{\formula{A}(t)}, just = \TAss
    [\sFmla{\True}{\eq[s][t]}, just = \TAss
      [\sFmla{\True}{\formula{A}(s)}, just = \TAss
        [\sFmla{\True}{\formula{A}(t)}, just={\TRule{\True}{\eq}[2, 3]}, close]
      ]
    ]
  ]
\end{oltableau}
Esse resultado talvez seja conhecido como o princípio da substituibilidade dos
idênticos, ou Lei de Leibniz.

!!^{tableau}s demonstram que $\eq$ é simétrica, isto é, que $\eq[s_1][s_2]
\Proves \eq[s_2][s_1]$:
\begin{oltableau}
  [\sFmla{\False}{\eq[s_2][s_1]}, just = \TAss
    [\sFmla{\True}{\eq[s_1][s_2]}, just = \TAss
      [\sFmla{\True}{\eq[s_1][s_1]}, just = {$\eq$}
        [\sFmla{\True}{\eq[s_2][s_1]}, just = {\TRule{\True}{\eq}[2, 3]}, close]
      ]
    ]
  ]
\end{oltableau}
Aqui, a linha $2$ é a primeira !!{formula} exigida,
$\sFmla{\True}{\eq[s_1][s_2]}$, por $\TRule{\True}{\eq}$.
A linha~$3$ é a segunda, da forma $\sFmla{\True}{!A(s_2)}$---considere
$!A(x)$ como
$\eq[x][s_1]$; então $!A(s_1)$ é $\eq[s_1][s_1]$ e $!A(s_2)$ é $\eq[s_2][s_1]$.

Eles também demonstram que $\eq$ é transitiva, isto é, que $\eq[s_1][s_2],
\eq[s_2][s_3] \Proves \eq[s_1][s_3]$:
\begin{oltableau}
  [\sFmla{\False}{\eq[s_1][s_3]}, just = \TAss
    [\sFmla{\True}{\eq[s_1][s_2]}, just = \TAss
      [\sFmla{\True}{\eq[s_2][s_3]}, just = \TAss
        [\sFmla{\True}{\eq[s_1][s_3]}, just = {\TRule{\True}{\eq}[3, 2]}, close]
      ]
    ]
  ]
\end{oltableau}
Nesse !!{tableau}, a primeira !!{formula} exigida por
$\TRule{\True}{\eq}$ é a linha~$3$, $\sFmla{\True}{\eq[s_2][s_3]}$
($s_2$ desempenha o papel de~$t_1$, e $s_3$ o papel de~$t_2$). A
segunda fórmula exigida, da forma~$\sFmla{\True}{!A(s_2)}$, é a linha~$2$.
Aqui, considere $!A(x)$ como $\eq[s_1][x]$; isso transforma $!A(s_2)$ em
$\eq[t_1][t_2]$ (isto é, a linha~$2$) e $!A(s_3)$ na
!!{formula}~$\eq[s_1][s_3]$ da conclusão.
\end{ex}

\begin{prob}
Forneça !!{tableau}s fechados para as seguintes fórmulas:
\begin{enumerate}
\item $\sFmla{\False}{\lforall[x][\lforall[y][((x = y \land !A(x))
      \lif !A(y))]]}$
\item $\sFmla{\False}{\lexists[x][(!A(x) \land
    \lforall[y][(!A(y) \lif y = x)])]}$,\\
  $\sFmla{\True}{\lexists[x][!A(x)] \land
    \lforall[y][\lforall[z][((!A(y) \land !A(z)) \lif y = z)]]}$
\end{enumerate}
\end{prob}

\end{document}
